(1) 将制冷机作为周期运行的封闭工质，其一个周期内能变化为零。热力学第一定律给 \(W_{\rm in}=Q_h-Q_c>0\)，其中机械功必须输入系统，才能把低温端热量搬运到更高温度的热库。 这里三个热功量均定义为正的大小，输入功不是热库自动输出的能量；因此正确的符号约定对于后续性能系数的计算尤其重要。
(2) 对可逆卡诺循环，总熵产生为零，两个热库的熵变之和满足 \(Q_h/T_h-Q_c/T_c=0\)，即 \(Q_h/Q_c=T_h/T_c\)。所以 \(W_{\rm in}=Q_c(T_h/T_c-1)\)，制冷性能系数 \(\mathrm{COP}_R=Q_c/W_{\rm in}=T_c/(T_h-T_c)\)。
(3) 当冷端制冷功率为 \(\dot Q_c\)，理想可逆机的最小输入功率为 \(\dot W_{\min}=\dot Q_c(T_h-T_c)/T_c\)，热端排热率为 \(\dot Q_h=\dot Q_c+\dot W_{\min}=\dot Q_cT_h/T_c\)。量纲均为瓦，显式满足能量守恒。
(4) 不可逆机器的热库熵产生非负，\(Q_h/T_h-Q_c/T_c\ge0\)，从而输入功不小于可逆值，COP不超过 \(T_c/(T_h-T_c)\)。当两热库温差趋零，理想COP趋大，但实际传热速率与功耗有限制；若固定热端温度、\(T_c\to0\)，COP趋零且所需功趋大。
